%Version 3.1 December 2024
% See section 11 of the User Manual for version history
%
%%%%%%%%%%%%%%%%%%%%%%%%%%%%%%%%%%%%%%%%%%%%%%%%%%%%%%%%%%%%%%%%%%%%%%
%%                                                                 %%
%% Please do not use \input{...} to include other tex files.       %%
%% Submit your LaTeX manuscript as one .tex document.              %%
%%                                                                 %%
%% All additional figures and files should be attached             %%
%% separately and not embedded in the \TeX\ document itself.       %%
%%                                                                 %%
%%%%%%%%%%%%%%%%%%%%%%%%%%%%%%%%%%%%%%%%%%%%%%%%%%%%%%%%%%%%%%%%%%%%%

%%\documentclass[referee,sn-basic]{sn-jnl}% referee option is meant for double line spacing

%%=======================================================%%
%% to print line numbers in the margin use lineno option %%
%%=======================================================%%

%%\documentclass[lineno,pdflatex,sn-basic]{sn-jnl}% Basic Springer Nature Reference Style/Chemistry Reference Style

%%=========================================================================================%%
%% the documentclass is set to pdflatex as default. You can delete it if not appropriate.  %%
%%=========================================================================================%%

%%\documentclass[sn-basic]{sn-jnl}% Basic Springer Nature Reference Style/Chemistry Reference Style

%%Note: the following reference styles support Namedate and Numbered referencing. By default the style follows the most common style. To switch between the options you can add or remove “Numbered” in the optional parenthesis. 
%%The option is available for: sn-basic.bst, sn-chicago.bst%  
 
% Toggle this for citation/layout mode:
%   \preprintmodetrue  -> author-year preprint citations like (Author et al., 2024)
%   \preprintmodefalse -> Nature submission formatting
% If you switch modes and natbib complains, run `latexmk -C sn-article.tex` once.
\newif\ifpreprintmode
\preprintmodetrue

\ifpreprintmode
\documentclass[pdflatex,sn-basic]{sn-jnl}% Preprint mode with author-year citations
\else
\documentclass[pdflatex,sn-nature]{sn-jnl}% Style for submissions to Nature Portfolio journals
\fi

% Toggle this for figure resolution:
%   \litefigurestrue  -> rasterized 400-DPI figures (figure_panels_lite/), ~5MB PDF, opens fast.
%   \litefiguresfalse -> original vector figures (figure_panels/), ~76MB, for final journal submission.
% Body text stays fully vector either way; only the data-dense figure panels are rasterized.
\newif\iflitefigures
\litefigurestrue
\iflitefigures
  \newcommand{\figdir}{figure_panels_lite}%
\else
  \newcommand{\figdir}{figure_panels}%
\fi
%%\documentclass[pdflatex,sn-basic]{sn-jnl}% Basic Springer Nature Reference Style/Chemistry Reference Style
%%\documentclass[pdflatex,sn-mathphys-num]{sn-jnl}% Math and Physical Sciences Numbered Reference Style
%%\documentclass[pdflatex,sn-mathphys-ay]{sn-jnl}% Math and Physical Sciences Author Year Reference Style
%%\documentclass[pdflatex,sn-aps]{sn-jnl}% American Physical Society (APS) Reference Style
%%\documentclass[pdflatex,sn-vancouver-num]{sn-jnl}% Vancouver Numbered Reference Style
%%\documentclass[pdflatex,sn-vancouver-ay]{sn-jnl}% Vancouver Author Year Reference Style
%%\documentclass[pdflatex,sn-apa]{sn-jnl}% APA Reference Style
%%\documentclass[pdflatex,sn-chicago]{sn-jnl}% Chicago-based Humanities Reference Style

%%%% Standard Packages
%%<additional latex packages if required can be included here>

\usepackage{graphicx}%
\usepackage{float}%
\usepackage{amsmath,amssymb,amsfonts}%
\usepackage{amsthm}%
\usepackage[title]{appendix}%
\usepackage{xcolor}%
\usepackage{enumitem}%
\usepackage{changepage}%

% --- Line numbers for the preprint only (Introduction onward) ---
\ifpreprintmode
\usepackage{lineno}%
% Make lineno number lines inside amsmath display environments too:
\newcommand*\patchAmsMathEnvironmentForLineno[1]{%
  \expandafter\let\csname old#1\expandafter\endcsname\csname #1\endcsname
  \expandafter\let\csname oldend#1\expandafter\endcsname\csname end#1\endcsname
  \renewenvironment{#1}%
    {\linenomath\csname old#1\endcsname}%
    {\csname oldend#1\endcsname\endlinenomath}}%
\newcommand*\patchBothAmsMathEnvironmentsForLineno[1]{%
  \patchAmsMathEnvironmentForLineno{#1}%
  \patchAmsMathEnvironmentForLineno{#1*}}%
\AtBeginDocument{%
  \patchBothAmsMathEnvironmentsForLineno{equation}%
  \patchBothAmsMathEnvironmentsForLineno{align}%
  \patchBothAmsMathEnvironmentsForLineno{flalign}%
  \patchBothAmsMathEnvironmentsForLineno{alignat}%
  \patchBothAmsMathEnvironmentsForLineno{gather}%
  \patchBothAmsMathEnvironmentsForLineno{multline}}%
\renewcommand\linenumberfont{\normalfont\footnotesize}%
\fi
%%%%

\newcommand{\figlegend}[1]{\par\smallskip\noindent\small #1\normalsize\par\medskip}

%%%%%=============================================================================%%%%
%%%%  Remarks: This template is provided to aid authors with the preparation
%%%%  of original research articles intended for submission to journals published 
%%%%  by Springer Nature. The guidance has been prepared in partnership with 
%%%%  production teams to conform to Springer Nature technical requirements. 
%%%%  Editorial and presentation requirements differ among journal portfolios and 
%%%%  research disciplines. You may find sections in this template are irrelevant 
%%%%  to your work and are empowered to omit any such section if allowed by the 
%%%%  journal you intend to submit to. The submission guidelines and policies 
%%%%  of the journal take precedence. A detailed User Manual is available in the 
%%%%  template package for technical guidance.
%%%%%=============================================================================%%%%

%% as per the requirement new theorem styles can be included as shown below
\theoremstyle{thmstyleone}%
\newtheorem{theorem}{Theorem}%  meant for continuous numbers
%%\newtheorem{theorem}{Theorem}[section]% meant for sectionwise numbers
%% optional argument [theorem] produces theorem numbering sequence instead of independent numbers for Proposition
\newtheorem{proposition}[theorem]{Proposition}% 
%%\newtheorem{proposition}{Proposition}% to get separate numbers for theorem and proposition etc.

\theoremstyle{thmstyletwo}%
\newtheorem{example}{Example}%
\newtheorem{remark}{Remark}%

\theoremstyle{thmstylethree}%
\newtheorem{definition}{Definition}%

\makeatletter
\def\Authorfont{\reset@font\fontsize{12bp}{14.5bp}\selectfont\boldmath\raggedright}%
\def\addressfont{\reset@font\fontsize{11bp}{13.5bp}\selectfont\raggedright}%
\let\SNorigprintabstract\printabstract
\let\SNorigprintkeywords\printkeywords
\renewcommand{\printabstract}{}
\renewcommand{\printkeywords}{}
\makeatother

\ifpreprintmode
\setcitestyle{authoryear,round,semicolon,aysep={,},notesep={; }}
\makeatletter
\def\hyper@natlinkbreak#1#2{#1}
\makeatother
\fi

\newcommand{\frontmatterstatement}[2]{%
  \par\medskip\noindent\textbf{#1}\par
  \noindent #2\par
}

\newcommand{\frontmatterblock}{%
  \frontmatterstatement{Funding statement}{\fundingstatementtext}
  \frontmatterstatement{Competing interests}{\competingintereststext}
  \frontmatterstatement{Acknowledgements}{\acknowledgementstext}
}

\newcommand{\fundingstatementtext}{This research was supported by the McNair Foundation and by NIH R01 MH129439, U01 NS121472, NINDS Research Education Grant Programs for Residents and Fellows in Neurology, Neurosurgery, Neuropathology, and Neuroradiology (UE5), the SNS Allan Friedman RUNN Research Grant, the NLM Training Program in Biomedical Informatics \& Data Science for Predoctoral \& Postdoctoral Fellows, T15LM007093-33, and the Gordon and Mary Cain Pediatric Neurology Research Foundation.}

\newcommand{\competingintereststext}{S.A.S. has consulting agreements with Boston Scientific, Zimmer Biomet, Koh Young, Abbott, and Neuropace. S.A.S. is co-founder of Motif Neurotech.}

\newcommand{\acknowledgementstext}{We thank Victoria Gates and Raissa Mathura for their assistance.}

\raggedbottom
%%\unnumbered% uncomment this for unnumbered level heads

\begin{document}

\title[Article Title]{Estimation of neuronal tuning for word meaning from passively recorded naturalistic speech}

%%=============================================================%%
%% GivenName	-> \fnm{Joergen W.}
%% Particle	-> \spfx{van der} -> surname prefix
%% FamilyName	-> \sur{Ploeg}
%% Suffix	-> \sfx{IV}
%% \author*[1,2]{\fnm{Joergen W.} \spfx{van der} \sur{Ploeg} 
%%  \sfx{IV}}\email{iauthor@gmail.com}
%%=============================================================%%

\author[1]{\fnm{Taha} \sur{Ismail}}

\author[1]{\fnm{Ana G.} \sur{Chavez}}

\author[1]{\fnm{Xinyuan} \sur{Yan}}

\author[1]{\fnm{Hanlin} \sur{Zhu}}

\author[1]{\fnm{Melissa} \sur{Franch}}

\author[1]{\fnm{James} \sur{Belanger}}

\author[1]{\fnm{Saipravallika} \sur{Chamarthi}}

\author[1]{\fnm{Katherine E.} \sur{Kabotyanski}}

\author[1]{\fnm{Kalman A.} \sur{Katlowitz}}

\author[1]{\fnm{Assia} \sur{Chericoni}}

\author[1]{\fnm{Elizabeth A.} \sur{Mickiewicz}}

\author[1]{\fnm{Timon} \sur{Merk}}

\author[1]{\fnm{Yewen} \sur{Zhou}}

\author[1]{\fnm{Nikhil D.} \sur{Shivakumar}}

\author[1]{\fnm{Paul J.} \sur{Steffan}}

\author[7]{\fnm{Rahul} \sur{Hingorani}}

\author[7]{\fnm{Mattson} \sur{Ogg}}

\author[7]{\fnm{Han} \sur{Yi}}

\author[1]{\fnm{Tomasz} \sur{Fraczek}}

\author[1,4]{\fnm{Eleonora} \sur{Bartoli}}

\author[2,5]{\fnm{Jay A.} \sur{Hennig}}

\author[1,2,3,4,5]{\fnm{Sameer A.} \sur{Sheth}}
\equalcont{These authors contributed equally to this work.}

\author[1,4,5,6]{\fnm{Nicole R.} \sur{Provenza}}
\equalcont{These authors contributed equally to this work.}

\author*[1,4,5]{\fnm{Benjamin Y.} \sur{Hayden}}\email{Benjamin.Hayden@bcm.edu}
\equalcont{These authors contributed equally to this work.}

\affil*[1]{\orgdiv{Department of Neurosurgery}, \orgname{Baylor College of Medicine}, \orgaddress{\city{Houston}, \state{TX}, \country{United States}}}

\affil[2]{\orgdiv{Department of Neuroscience}, \orgname{Baylor College of Medicine}, \orgaddress{\city{Houston}, \state{TX}, \country{United States}}}

\affil[3]{\orgdiv{Department of Psychiatry and Behavioral Sciences}, \orgname{Baylor College of Medicine}, \orgaddress{\city{Houston}, \state{TX}, \country{United States}}}

\affil[4]{\orgdiv{Department of Electrical and Computer Engineering}, \orgname{Rice University}, \orgaddress{\city{Houston}, \state{TX}, \country{United States}}}

\affil[5]{\orgdiv{Neuroengineering Initiative}, \orgname{Rice University}, \orgaddress{\city{Houston}, \state{TX}, \country{United States}}}

\affil[6]{\orgdiv{Department of Bioengineering}, \orgname{Rice University}, \orgaddress{\city{Houston}, \state{TX}, \country{United States}}}

\affil[7]{\orgdiv{Research and Exploratory Development Department}, \orgname{Johns Hopkins University Applied Physics Laboratory}, \orgaddress{\city{Laurel}, \state{MD}, \country{United States}}}

%%==================================%%
%% Sample for unstructured abstract %%
%%==================================%%

\abstract{The ability to derive neural-level language coding models holds great scientific and clinical potential. Current approaches are limited by the scale and ethological validity of input data; applications requiring large, rare, or naturalistic samples in particular would benefit from the ability to infer neural coding from incidental everyday speech. Here we present a novel pipeline designed to leverage spontaneous and incidental naturalistic speech. This pipeline performs transcription, segmentation, and video-assisted diarization, as well as alignment and spike detection of neural data. We apply this pipeline to a dataset derived from 21 patients (6+ days each, over 800 hours and 5 million words total). We benchmark both encoding and decoding models against extensive and rare ground-truth control datasets consisting of human-curated word-level temporal alignment and manually sorted spikes. We further validate our approach by quantifying representational drift, effect of dataset size, and differences between six brain areas. Together, these findings demonstrate that incidental natural speech is sufficiently processed in the brain to enable the estimation of neural-level embeddings.}

%%================================%%
%% Sample for structured abstract %%
%%================================%%

% \abstract{\textbf{Purpose:} The abstract serves both as a general introduction to the topic and as a brief, non-technical summary of the main results and their implications. The abstract must not include subheadings (unless expressly permitted in the journal's Instructions to Authors), equations or citations. As a guide the abstract should not exceed 200 words. Most journals do not set a hard limit however authors are advised to check the author instructions for the journal they are submitting to.
% 
% \textbf{Methods:} The abstract serves both as a general introduction to the topic and as a brief, non-technical summary of the main results and their implications. The abstract must not include subheadings (unless expressly permitted in the journal's Instructions to Authors), equations or citations. As a guide the abstract should not exceed 200 words. Most journals do not set a hard limit however authors are advised to check the author instructions for the journal they are submitting to.
% 
% \textbf{Results:} The abstract serves both as a general introduction to the topic and as a brief, non-technical summary of the main results and their implications. The abstract must not include subheadings (unless expressly permitted in the journal's Instructions to Authors), equations or citations. As a guide the abstract should not exceed 200 words. Most journals do not set a hard limit however authors are advised to check the author instructions for the journal they are submitting to.
% 
% \textbf{Conclusion:} The abstract serves both as a general introduction to the topic and as a brief, non-technical summary of the main results and their implications. The abstract must not include subheadings (unless expressly permitted in the journal's Instructions to Authors), equations or citations. As a guide the abstract should not exceed 200 words. Most journals do not set a hard limit however authors are advised to check the author instructions for the journal they are submitting to.}

\keywords{semantics, neural, encoding, decoding, continuous}

%%\pacs[JEL Classification]{D8, H51}

%%\pacs[MSC Classification]{35A01, 65L10, 65L12, 65L20, 65L70}

\makeatletter
\gdef\@presentaddresstext{\frontmatterblock}
\global\presentaddresstrue
\makeatother

\newgeometry{top=0.8in,bottom=0.8in,left=0.9in,right=0.9in}
\maketitle

\newgeometry{margin=1in}
\clearpage
\makeatletter
\SNorigprintabstract
\SNorigprintkeywords
\makeatother
\clearpage
\ifpreprintmode\linenumbers\fi
\section{Introduction}\label{sec1}

As we listen to speech, our brains continuously transform that input into
structured neural representations. By combining measurements of
neural responses with embeddings derived from large-language models, scientists
can quantify the association between neural responses and the language that
evokes them. This approach works with a range of neural measures, including
fMRI, MEG, EEG, and intracranial EEG \citep{bib1,bib2,bib3,bib4,bib5,bib9,bib10,bib11}.
While each of these measures have their value, neuron-level recordings track both
the high temporal resolution of speech and give the high spatial resolution of
neural computation \citep{bib13,bib14,bib4,bib5,bib12,bib48,bib49}. As a result, they are
often used in brain--computer interfaces (BCIs; \citealp{bib102,bib16,bib17}), as well as in
basic science studies focused on mechanism \citep[e.g.,][]{bib134}.

Neuron-level language representations can be formalized using \textit{encoding
models}, which relate patterns of brain activity to language model-derived
embeddings \citep{bib2,bib19,bib10,bib20,bib4,bib5,bib21,bib22}. \textit{Decoding
models} provide a complementary perspective by asking whether neural data can
classify linguistic stimuli \citep{bib23,bib2,bib24,bib25}. Both have potential
value for both clinical and basic science; moreover, their joint application
potentially allows for both improved interpretability and enhanced understanding \citep{bib68,bib135}.

Most such studies rely on brief, highly controlled laboratory paradigms with
laborious manual curation of speech (transcription, segmentation, diarization)
and manual processing of neural data. As a
result, it remains unknown whether neural-level robust semantic models can be recovered from
long-duration, naturalistic speech encountered in everyday life, and, if so, how
accurate they will be \citep{bib27,bib28,bib29,bib30,bib31}. Modeling neural responses from incidental
continuously recorded speech introduces several challenges
\citep{bib32,bib27,bib22}. Everyday speech contains frequent dysfluencies,
variable articulation, speaker overlap, and shifting acoustic contexts, all of
which complicate precise temporal alignment between language and neural activity
\citep{bib33,bib34,bib35,bib36}. These properties make accurate transcription,
diarization, and segmentation difficult. Moreover, the long durations of natural
conversation preclude gold-standard manual spike sorting across long recordings
\citep{bib37,bib38}. Although emerging automated pipelines for speech processing
and spike detection now make approximate spike sorting in such datasets
tractable, these tools remain imperfect and can introduce additional sources of
noise and bias \citep{bib39,bib105,bib106,bib107,bib108,bib109,bib110,bib111}. One of the biggest barriers to evaluating these
models is the lack of high-quality ground-truthed datasets large enough to
evaluate automated pipelines. Despite these barriers, we hypothesize that the brain processes natural speech well enough that automated pipelines can reliably be deployed to recover valid speech models.

In any case, naturalistic speech offers unique scientific and
practical advantages. Daily life provides an exceptionally rich and diverse
linguistic input---often tens of thousands of words per day---far exceeding what
can be presented in conventional experiments \citep{bib32,bib22,bib40,bib27}.
This scale is critical for studying rare linguistic events, context-dependent
meaning, and phenomena such as dysfluencies that can be difficult to elicit under
controlled conditions. In addition, continuously updated models derived from
everyday speech are increasingly relevant for clinical technologies, including
speech prosthetics and adaptive deep brain stimulation, which require robust
neural--linguistic mappings outside the laboratory \citep{bib102,bib41,bib42}.
Here we describe a novel pipeline deployed to analyze all incidental speech from
within the EMU at a single hospital for 21 patients over more than 3000 of total monitoring time. Critically,
we leverage a ground-truth dataset derived from real-world conversations that is
manually aligned and diarized, along with corresponding manual spike sorting.

\clearpage
\section{Results}\label{sec2}

We recorded speech and brain activity over a period of 6 to 11 days for 21 adult patients with medically refractory epilepsy undergoing intracranial seizure monitoring. We continuously recorded speech occurring during the entirety of each patient's stay (average: 6.72 $\pm$ 0.52 days) in our epilepsy monitoring unit (EMU, \textbf{Methods}). Two studio microphones, one directly above the patient's bed and another to the side, recorded 24/7 speech occurring within the patient room. Additionally, eight cameras located throughout the room, capturing various angles of the patient bed, recorded patient activity during their stay (\textbf{Figure 1A}). The microphones and cameras recorded many conversations with friends, family, clinical and research staff, as well as speech played from the television (\textbf{Figure 1B}). Intracranial EEG electrodes with microwires for collecting single-neuron spiking data (median: 56 contacts per patient, ranging from 16--96 contacts in various brain regions) recorded single and multi-unit spiking activity at 30 kHz for the entire patient stay, allowing us to sync microwire unit data at high temporal resolution.

All speech data were processed through a voice activity detection model \citep{bib43} to find passages of audio data that contained identifiable human speech. Speech was then segmented into sentences and words using WhisperX \citep{bib39}. Multiple scoring metrics, including CTC log-likelihood, WhisperX quality score, spectral entropy, and DNSMOS quality scores were calculated for each sentence (\textbf{Methods}). Sentences at low-quality extrema (outside 0.75 $\times$ IQR) of their relative distributions were removed from the dataset (\textbf{Figure 1C}). We used a standard autothresholding algorithm \citep{bib112,bib113,bib114} to capture multi-unit spiking activity that could be aligned with word timestamps (\textbf{Figure 1C}, \textbf{Methods}). For convenience, we refer to each of these autolabelled sets as \textit{units}.

To identify segments of patient speech, we used a multimodal active speaker detection (ASD) model \citep{bib44}. Audio-only diarization methods make frequent errors in naturalistic recording conditions, such as overlapping speakers and background noise, which led to us turning to multimodal methods \citep{bib45,bib46}. Specifically, we combined video with audio features to detect facial movements associated with speech to identify video frames in which the patient is likely to be speaking (\textbf{Figure 1D}, \textbf{Methods}). We used the speaker labels from this model to compare patient-only versus all speech in our downstream semantic coding models.

Our \textbf{encoding models} involved regressing GPT-2 embeddings against spike counts for individual units \citep[cf.][]{bib4,bib48,bib49}. The model was implemented as a Poisson GLM with an L2 regularizer and a duration term included as an offset variable (\textbf{Figure 1E}). Models were fitted independently across neurons in bulk matrices in PyTorch (\textbf{Methods}). GPT-2 embeddings were calculated by taking the hidden state of the last layer of GPT-2-large from an inference pass of the current word, with the prior 200 words in the continuous speech stream taken as context (\textbf{Methods}). The final hidden layer was selected due to it being the natural contextual representation of a word as it is directly used for next-token prediction in the transformer architecture \citep{bib47}. Moreover, prior studies find qualitatively robust performance for neural encoders in the final layer of GPT-2 as well as earlier layers \citep{bib48,bib4,bib49}.

Our \textbf{decoding models} involved training an XGBoost classifier with population firing rate activity (defined as the average neural firing rate across a word interval for a set of neurons) to predict broad semantic category labels, which included classes such as ``relationships'', ``numerical'', and ``medical'' (\textbf{Figure 1F}). XGBoost was chosen for its skill in learning nonlinear classifier boundaries, which is helpful for neural decoding \citep{bib50}. Moreover, XGBoost is frequently used with success for behavioral decoding from invasive neural recordings \citep{bib51,bib52}.

\clearpage

\begin{figure}[H]
\centering
\includegraphics[width=\textwidth,height=0.86\textheight,keepaspectratio]{\figdir/Figure_1_FINAL_Jun_16.pdf}
\label{fig:figure1}
\end{figure}
\figlegend{\textbf{Figure 1 -- Neural-behavioral sync and modeling pipeline.} \textbf{A.} Picture and diagram showing the Epilepsy Monitoring Unit (EMU) where patients are kept for intracranial seizure monitoring, with icons displaying the audio and video recording modalities present in the room and for research use. \textbf{B.} Synchronization of audio and neural signal courtesy of our Blackrock Neurotech recording system that maintains a synchronized stream of both, sampled at 30 kHz. \textbf{C.} Automated speech extraction and autothresholding pipeline. \textbf{D.} Methodology for extracting patient speaker scores from the audio-video stream and filtering patient-only speech segments. \textbf{E.} Input structure of our Poisson GLM used for semantic and emotion encoding. \textbf{F.} Input and model structure used for XGBoost semantic category decoding.}

\clearpage
\subsection{Dataset spanning millions of words with accurate transcription and diarization}

We recorded a total of 5{,}270{,}665 words both spoken and heard by 21 participants with simultaneous brain recordings. The number of words for each patient ranged from 40{,}351 to 521{,}360. Seventeen of the patients had datasets with over 100{,}000 words. Across all patients, we collected a total of 871.2 hours of speech (spoken and/or heard), ranging between 7.1 and over 87.7 hours of speech per patient (\textbf{Figure 2A}). Across the fourteen patients for which we have quality video data available, we collected a total of 156{,}397 words of patient-specific speech, ranging from 1{,}644 to 33{,}972 words per patient. We find the lexicon of spoken words to span a range of expected topics, including a large mode for medical terminology (\textbf{Figure 2B}). Additionally, we find that the timing of speech production and listening follows a cyclic pattern, with the number of words spoken peaking in the early afternoon hours and declining significantly late at night (\textbf{Figure 2C}).

\textit{Ground-truth accuracy:} We first quantified the accuracy of the WhisperX transcription and video-based diarization for our specific EMU-based context. For a subset of fifteen patients, we used sections of data in which study staff had planned conversations with the patient ($\sim$45 minutes/patient). We then performed precise manual temporal alignment and speaker identification (i.e., diarization) of every word spoken and heard \citep{bib53}. We used these data as a ground truth to evaluate the accuracy of the automated transcription and diarization of the same conversation. For transcription accuracy, we evaluated every sentence in our manual transcript with its corresponding match in the automatic transcript. We focused on semantically meaningful sentences, which included sentences with at least two non-filler words (filler words here refers to words like ``\textit{um}'', ``\textit{hmm}'', and ``\textit{haha}''). For each sentence in the manual transcript we searched for its corresponding sentence (if any) in the automated transcript within a 60 second window around the sentence. A match was considered to be the sentence with the lowest word error rate (WER) in the window as long as it was 0.5 or less, which was determined via Levenshtein distances \citep{bib115} (\textbf{Methods}). We found that, on average, 90.15\% of semantically meaningful sentences were matched in the automated transcript, with the automated transcript capturing the majority of sentences in the conversation for all patients (\textbf{Figure 2D}). This correspondence indicates WhisperX is generally successful in capturing instances of speech in the EMU, with infrequent errors resulting from known variables such as overlapping speakers and noisy audio segments \citep{bib39}.

\textit{Quantifying semantic precision:} We calculated four summary statistics for sentence matches in order to assess the extent of transcription errors. The average word error rate similarity score was 0.89, indicating that only 11\% of the characters in automated transcripts differed from the true transcript on average. We next calculated the cosine similarity between matched sentence embeddings (derived via the all-mpnet-base-v2 version of the SBERT sentence transformers model; \citealp{bib54}). We found the mean cosine similarity score to also be 0.89. When looking at non-perfect matches (i.e., WER similarity $<$1.0), the average SBERT cosine similarity score is greater than the average WER similarity score ($p<0.001$). As such, even in the cases where minor word-level or syntactical errors are made, semantic meaning tended to be maintained for sentence matches. For example, the sentences ``I knew he had taken it \textit{because} I was just like tutoring with him \textit{one on one}'' and ``I knew he had taken it \textit{cause} I was just like tutoring with him,'' had a slightly lower WER similarity score of 0.8 whereas their SBERT cosine similarity is 0.981 (\textbf{Figure 2E}).

\textit{Quantifying temporal precision:} We also calculated the timestamp error of all matched sentences and words within those sentences to quantify temporal accuracy in word segmentation. Notably, the median absolute distance between matched sentences (defined from the midpoints of the sentences) was 115.50 milliseconds, and the median absolute distance between words within each matched sentence was only 71.68 milliseconds. This minimal discrepancy indicates timestamp errors are relatively minor and should not significantly affect estimates of semantic tuning. Overall, we found that transcription and segmentation errors were relatively minor and should not invalidate the possibility of generating semantic coding models from naturalistic speech (\textbf{Figure 2F}).

To assess our diarization accuracy, we compared the outputs of LR-ASD to the manually annotated speakers for each sentence. LR-ASD provides a speaker confidence score for every frame of a video, so we labeled a speech segment as coming from the patient if the model labeled the patient as speaking for the majority of frames. Labelling in this way notably struggles in sections with overlapping speakers, similar to automated transcription. When calculating basic classification metrics and comparing them to each patient's base speaking rate (i.e., the portion of the conversation they spent talking), we find that the video-diarization model is significantly more precise than the base speaking rate for all tested patients for which we had video data available (\textbf{Figure 2G}). Moreover, for 50\% of patients with video available, LR-ASD is more than 70\% accurate when classifying a sentence as belonging to the patient. This result is in line with other groups' findings that multimodal diarization models significantly outperform audio-only methods \citep{bib55,bib56}. We find that our pipeline struggles with recall for many patients dependent on video quality, meaning that while we can be reasonably confident that a sentence marked as ``patient'' is patient speech, we cannot be confident that we are necessarily detecting the majority of patient-spoken sentences. As such, for downstream analyses we do not make the assumption that sentences not classified as patient speech necessarily are spoken by other people. Overall, while automated diarization is not as comprehensively accurate as automated transcription, we consider it to be sufficiently reliable for analyses evaluating semantic tuning of patient-specific speech.

\clearpage

\begin{figure}[H]
\centering
\includegraphics[width=\textwidth,height=0.86\textheight,keepaspectratio]{\figdir/Figure_2_Jun23.pdf}
\label{fig:figure2}
\end{figure}
\figlegend{\textbf{Figure 2 -- Data quantity and annotation accuracy.} \textbf{A.} Bar plots showing the amount of speech data collected from our pipeline per patient, sorted by number of words. \textbf{B.} 2D UMAP plot of all words collected from our pipeline, clustered by word2vec embeddings and colored by semantic category (patient/vistor named redacted for compliance). Words generally cluster by their assigned semantic category. \textbf{C.} Line plot showing the average number of words across patients collected per 3 hour block starting the first midnight after patient admission, displaying a daily cyclic pattern. Background color displays the standard error per time block across patients. \textbf{D.} Lollipop plot displaying the portion of conversation sentences for which a corresponding match was found in the associated automated transcript. \textbf{E.} Example sentences displaying what potential pairs of matched sentences looked like and the types of transcription errors corrected for. \textbf{F.} Histograms describing transcription and segmentation errors across matched sentences in terms of WER similarity, SBERT cosine similarity, sentence timing errors, and word timing errors. \textbf{G.} Paired barplot summarizing diarization accuracy of annotated conversations compared to patient base speaking rate.}

\clearpage
\subsection{Significant semantic encoding in the continuous speech dataset}

We first aimed to understand whether our automated pipeline allows for the development of meaningful semantic encoding models. We analyzed recovered autothresholded unit data from all available brain structures: (1) hippocampus (HPC, 496 units), (2) amygdala (AMY, 88 units), (3) orbitofrontal cortex (OFC, 104 units), (4) thalamus (THAL, centromedian nucleus, 64 units), (5) anterior cingulate cortex (ACC, 240 units), and (6) posterior cingulate cortex (PCC, 24 units). We first examined all autothresholded units together. We collected between 16 to 96 units per patient, with the most units coming from the hippocampus and anterior cingulate and the least coming from the thalamus and posterior cingulate.

We evaluated \textit{encoder model performance} using the pseudo R\textsuperscript{{2}} score, which defines the relative performance gain over a null model that predicts spike counts from the average firing rate and duration of each word (\textbf{Methods}). Pseudo R\textsuperscript{{2}} was chosen for its interpretability and inherent benchmarking: a score equal to zero indicates performance that is equal to our specified null model. Negative values indicate performance that was worse than the null model. Moreover, pseudo R\textsuperscript{{2}} scores with mean firing rate nulls are commonly used to evaluate GLM performance in neural coding analyses \citep{bib57,bib58,bib59}. Our automated pipeline produced robust encoding models, as seen from average pseudo R\textsuperscript{{2}} scores above 0 ($0.004 \pm 0.0003$ to $0.12 \pm 0.025$) for all patients (\textbf{Figure 3A}). Our response range is typical for neural encoding models \citep{bib57,bib58,bib59}, and significant scores above zero confirm that our automated pipeline is able to detect semantic tuning without manual data curation.

We next compared encoding performance for patient-only speech identified by our video diarization pipeline with all ambient speech. We overall found encoding performance to be stronger for produced speech than all speech (\textbf{Figure 3B}). Specifically, we see a 2.42-fold increase (from 0.026 to 0.064) in pseudo R\textsuperscript{{2}} scores for patient-spoken speech. This finding is especially notable considering our patient-only models were trained on $50.52~(\pm 12.30)$ times less data on average. One reason produced speech may be more strongly encoded is that it requires attention, while ambient speech may be ignored \citep{bib35}.

\subsubsection{Comparison with controlled tasks}

We next compared embeddings derived from incidental speech to those derived from two highly controlled tasks: (1) an unconstrained conversation task \citep{bib53,bib49}, and (2) a podcast listening task \citep{bib4,bib48}; (\textbf{Methods}). We found that encoding strength is significantly higher for our controlled datasets, with a $4.43~(\pm 1.67)$-fold average greater encoding strength for the conversation task and an $8.59~(\pm 3.12)$-fold greater encoding strength for the podcast task (\textbf{Figure 3C}).

To gain insight into why the controlled tasks produce stronger encoding, we replicated this analysis using a version of both controlled datasets with autothresholded spikes instead of manually curated spikes. We found that using autothresholded spikes significantly improves encoder performance, with a $6.62~(\pm 1.58)$-fold average increase for the conversation task and a $1.47~(\pm 0.12)$-fold increase for the podcast task. Prior work has demonstrated that multiunit activity performs comparably to sorted spikes for population coding tasks and may even filter out useful signal \citep{bib60,bib61}, which may explain why spike sorting does not improve our model performance in the aggregate.

We then tested a version using uncorrected WhisperX transcripts rather than hand-labelled transcripts. We observed no significant difference in performance from using uncorrected WhisperX transcripts compared to hand-labelled ones, with a $0.95~(\pm 0.037)$-fold change in the conversation task and $1.01~(\pm 0.01)$-fold change in the podcast task. This result suggests manual annotation may have modest, if any, benefits for derivation of semantic encoding models. Together, these results suggest that poorer performance in the naturalistic dataset likely occurs for reasons other than the limitations of automated data processing---such as representational drift or patient attention variability.

\subsubsection{Variability in encoding over time}

We next evaluated how pooling data across days affects our model performance compared to training models on single-day data. We trained encoding models from each individual day in each patient's stay and compared them to models trained on the entire patient stay (averaging 6.72 days, ranging from 6--11 days). We find that models trained on individual days performed significantly better than models trained with data from all days, with $3.06~(\pm 0.44)$-fold greater pseudo R\textsuperscript{{2}} scores on average (\textbf{Figure 3D}). Thus, day-to-day variability is a major driver of variation in encoding model performance.

While during formal experiments we are able to limit the number of distractions and ensure patients are focused on the task, directing patient attention is not possible in a naturalistic setting, and neural tuning to semantics may differ based on the stimulus the patient is (or is not) paying attention to \citep{bib35}. To probe attentional changes in model performance, we trained encoding models on two-hour epochs across the patient stay, such that we could compare encoding strength at 2 PM and 7 PM on the same day, for example. We assessed whether epoch encoding strength correlated with the portion of time in the epoch the patient spent speaking, which can potentially be viewed as a proxy for patient attention and engagement with ambient speech. For example, the portion would be 0 when the patient is passively listening to the television while engaged in a separate activity, while it could be as high as 0.7 when they are actively engaged in conversation with another person. We find that, when regressing patient speech fraction against epoch encoding strength via weighted least squares (see: \textbf{Methods}), 62\% of patients have a positive trend (56\% significant), while 38\% of patients have a negative trend (19\% significant, \textbf{Figure 3E}). While this result may be a consequence of our units being significantly more tuned to patient speech, it may also indicate that overall, semantic tuning in our units is significantly correlated with the portion of patient speech engagement, potentially indicating that patient engagement and/or attention is a significant factor in developing high-performing semantic coding models.

\subsubsection{Comparing brain regions}

We found all recorded brain regions (median pseudo R\textsuperscript{{2}} $>$ 0) generate significant semantic encoding models from the naturalistic speech dataset. We found PCC, ACC and HPC to be the highest performing brain regions, with the remaining brain regions ranked in the following decreasing order: thalamus, OFC, and amygdala (\textbf{Figure 3F}). This result corroborates findings from existing literature implicating hippocampus in higher-level semantic/symbolic language processing and findings of distributed semantic networks throughout the brain \citep[e.g.,][]{bib2}.

\clearpage

\begin{figure}[H]
\centering
\includegraphics[width=\textwidth,height=0.86\textheight,keepaspectratio]{\figdir/Figure_3_Jun16.pdf}
\label{fig:figure3}
\end{figure}
\figlegend{\textbf{Figure 3 -- Significant semantic coding from naturalistic speech data.} \textbf{A.} Bar plot showing average semantic encoding performance per patient across all significant units with standard error bars, sorted by model performance. Number of significant units are listed above each bar. \textbf{B.} Paired line plot comparing semantic encoding performance between all speech and patient-only speech conditions. \textbf{C.} Grouped bar plot showing comparative semantic encoding performance between the naturalistic speech and experimental datasets (conversation + podcast), with automated conditions for each task. \textbf{D.} Paired line plot comparing semantic encoding performance between all speech and within-day speech conditions. \textbf{E.} Line plot with best fit lines and summary charts showing relationship between epoch encoding strength and portion of time patient spent speaking. \textbf{F.} Boxplot showing semantic encoding model performance distribution by brain region.}

\clearpage
\subsection{Significant semantic decoding in the continuous speech dataset}

Decoding models of semantics play a complementary role to encoding models and are especially interesting for BCI purposes \citep{bib16,bib62,bib22,bib63}. To simplify the semantic decoding problem, we classified all spoken words into one of ten semantic categories \citep[\textbf{Methods}; cf.][]{bib4}. We evaluated decoder model performance using overall accuracy scores (ratio of correctly identified words to total number of words). With ten semantic categories, chance performance was 10\%.

Our pipeline produces significant decoding. Models for all patients but one decode semantics with more than 20\% accuracy (average accuracy: $20.9\% \pm 0.08\%$, \textbf{Figure 4A}). In contrast to our encoding models, we find that decoding is more accurate (although modestly so) for all speech versus patient-only speech. Specifically, we see $18\%~(\pm 2.6\%)$ lower classification accuracy for patient-only speech (\textbf{Figure 4B}).

\subsubsection{Decoding: comparison with ground-truth datasets}

Unlike encoding, decoding models trained on all speech perform better than models trained in our ground-truth datasets: we saw a $36\%~(\pm 1.8\%)$ lower performance for the conversation task and an $11\%~(\pm 2.5\%)$ lower performance for the podcast task (\textbf{Figure 4C}). Decoding may be more robust to confounds (patient attention, variable neural quality, etc.) not found in our controlled tasks than our single-unit confounds are. As a result, population decoding models may benefit more from the additional training data available via automated parsing from naturalistic environments.

As with our encoding analyses, manual curation of spikes and transcript labels do not significantly improve the accuracy of semantic decoder models for our controlled tasks. We find no significant difference (i.e., overlapping standard errors) in the decoder accuracy for models trained with sorted versus autothresholded spikes. Moreover, while we find decoder performance is $4.0\%~(\pm 1.3\%)$ less accurate with automated transcripts in the podcast task, decoder performance is actually slightly more accurate ($18.3\% \pm 3.2\%$) with automated transcripts in the conversation task. Overall, decoders appear to be more robust to the inaccuracies produced by automated neural and behavioral processing models than encoding models.

\subsubsection{Decoding: variability in performance over time}

Decoding models are more stable over time than encoding models. Indeed, single-day model performance is $9\%~(3.5\%)$ worse than full multi-day dataset decoding (\textbf{Figure 4D}). Thus, the benefit of additional training data outweighs the impediment of cross-day neural variability.

We also find that decoding accuracy is better for patients with more recorded units than fewer (\textbf{Figure 4E}). This result provides further support for the notion that larger populations of neural data allow for the training of more accurate models of human behavior that take advantage of the redundancy of single-neuron tuning \citep{bib64,bib65,bib61} and are robust to failures and fluctuations that occur at the single-unit level.

We also assessed the variability of semantic decoding across the two-hour epoch scale and probed whether the same attentional variation as in encoding emerged. Overall, we find a similar if less pronounced trend, where 53\% of patients have a positive trend (20\% significant), whilst 47\% of patients have a negative trend (6\% significant, \textbf{Figure 4F}). This result mirrors our encoding findings and corroborates the idea that patient engagement may improve the performance of our semantic coding models. Decoding performance is also more stable across 2-hour epochs than encoding performance: our epoch decoding performance has an $85.9\%~(\pm 0.16\%)$ smaller coefficient of variation than our epoch encoding performance ($p<0.0001$; Feltz--Miller test, see: \textbf{Methods}, \textbf{Figure 4G}).

\subsubsection{Decoding: comparing brain regions}

We trained decoding models separately for each brain region. All brain regions performed significantly above chance, but worse than models trained using all available units across regions ($p<0.001$). As with the encoding models (see above), the PCC and HPC are the highest performing regions, followed by the thalamus and ACC (\textbf{Figure 4H}).

\clearpage

\begin{figure}[H]
\centering
\includegraphics[width=\textwidth,height=0.86\textheight,keepaspectratio]{\figdir/Figure_4_Jun16.pdf}
\label{fig:figure4}
\end{figure}
\figlegend{\textbf{Figure 4 -- Significant semantic decoding from naturalistic speech data.} \textbf{A.} Bar plot showing average semantic decoding performance per patient across all significant units with standard error bars, sorted by model performance. Number of significant models are listed above each bar. \textbf{B.} Paired line plot comparing semantic decoding performance between all speech and patient-only speech conditions. \textbf{C.} Grouped bar plot showing comparative semantic decoding performance between the naturalistic speech and experimental datasets (conversation + podcast), with automated conditions for each task. \textbf{D.} Paired line plot comparing semantic decoding performance between all speech and within-day speech conditions. \textbf{E.} Scatterplot with line of best fit highlighting trend between unit population size and model performance. \textbf{F.} Line plot with best fit lines and summary charts showing relationship between epoch decoding strength and portion of time patient spent speaking. \textbf{G.} Line plots showcasing epoch model performance variation over time. \textbf{H.} Boxplot showing semantic decoding model performance distribution by brain region.}

\clearpage
\subsection{Encoding model performance as a function of dataset size}

The scale of our dataset lets us determine the effect of data quantity on performance. When we randomly decimate our per-day datasets to portions ranging from 1\% to 90\% of the data (\textbf{Methods}), we generally find that additional training data improves encoding model performance (\textbf{Figure 5A}). That performance shows marginal decline, although we do not detect any saturation. Our best fit curve is modelled after the Hill equation, which is commonly used for sigmoidal biological relationships (\textbf{Methods}). Moreover, we found the Hill equation to significantly fit our data according to permutation testing ($p<0.01$) and it also better explained experimental data than a weighted least squares fit ($\mathrm{Loss}_{\mathrm{hill}} < \mathrm{Loss}_{\mathrm{wls}}$; $d\mathrm{Loss} = 90.4$).

We next quantified the minimum amount of data needed to generate statistically meaningful results. We calculate pseudo R\textsuperscript{{2}} by comparing it to a simpler model using word duration only. We find that model performance exceeds durational null performance when using at least $4.7\%~(\pm 0.04\%)$ of the dataset on average.

We then repeated these analyses using number of words, instead of number of minutes. As with minutes, we find a clear logistic growth curve, with performance exceeding null at around $724~(\pm 3)$ words (\textbf{Figure 5B}). The Hill equation similarly significantly fit our data and performed better than a simple linear fit ($\mathrm{Loss}_{\mathrm{hill}} < \mathrm{Loss}_{\mathrm{wls}}$; $d\mathrm{Loss} = -307.5$).

We find that the proportion of units significantly tuned to our semantic embeddings also increased according to a logistic growth curve, with 1\% of data yielding an average of $22.2\%~(\pm 0.2\%)$ significant units before saturating to an average of $91.6\%~(\pm 0.6\%)$ of units at around 70\% of each dataset (\textbf{Figure 5C}). Contextual transformer embeddings capture many aspects of language relevant to word meaning---including syntax, grammar, word order, surrounding context---within a distributed high-dimensional representation \citep{bib66}, which means a unit significantly tuned to these embeddings could be responsive to any combination of these features. Moreover, the general mixed selectivity of neurons in higher cortical regions \citep{bib67} may explain how weak correlations with some dimensions of the transformer embedding space may reveal themselves at sufficient sample size. Our analysis showcases a classic power curve, in that only the most prominent effects are recognizable at smaller sample sizes whereas weaker effects only emerge when there is enough data to detect them.

\clearpage

\begin{figure}[H]
\centering
\includegraphics[width=\textwidth,height=0.86\textheight,keepaspectratio]{\figdir/Figure_5_Jun12.pdf}
\label{fig:figure5}
\end{figure}
\figlegend{\textbf{Figure 5 -- Decimation analysis.} \textbf{A.} Boxplot with Hill curve showing semantic encoding performance growth as function of dataset portion used for training. \textbf{B.} 2D histogram with Hill curve showing semantic encoding performance growth as function of words in sample. \textbf{C.} Scatterplot of fraction of significant semantic encoding units as a function of dataset portions, with portion averages as an overlaid line plot.}

\clearpage
\subsection{Assessing differences in coding models across time (functional drift)}

We wanted to assess the stability of our encoding and decoding model performance due to known phenomena of representational drift and recording inconsistencies \citep{bib136,bib137}. Importantly, we wanted to address the rate of model degradation over time and whether it differed for unit encoding versus population decoding models, as such questions hold significant ramifications for the development of generalizable neural coding models of behavior.

We find that, for encoding models, model performance decreases as a function of the temporal separation between the original training day and the subsequent test day ($p<0.001$, \textbf{Figure 6A}). This effect is strongest when comparing within-day testing to the next-day ($+1$), where the median pseudo R\textsuperscript{{2}} score dips from $0.034~(\pm 0.006)$ to $-0.027~(\pm 0.008)$ and tends to grow more negative over time ($R^2$ scores are normalized in reference to null performance, so decreases below zero are meaningful). Moreover, the portion of significant units plummets from $97.8\%~(\pm 1.2\%)$ on average to $7.8\%~(2.4\%)$ after just one day of separation (\textbf{Figure 6B}). Overall, we find that significant model degradation occurs for our single-unit encoding models after even just one day of temporal separation, rendering most models unreliable for downstream use cases after their original training date.

Unlike encoding models, decoding model performance does not significantly decrease with distance in time ($p = 0.6450$, \textbf{Figure 6C}). However, similar to encoding, the proportion of significant models does decrease significantly over time (weighted least squares, $p<0.001$), albeit at a much more steady rate (\textbf{Figure 6D}). Specifically, the percent of significant models starts at $95.3\%~(\pm 1.9\%)$ and degrades to $68.6\%~(\pm 9.5\%)$ of models still performing with statistical significance after eight days from their original training date. The relative stability of decoding drift is apparent when we compare the distributions of regression slopes of individual model performance between encoding and decoding. We found model decoding slopes to be significantly higher than model encoding slopes (one-tailed $t$-test, $p<0.001$, \textbf{Figure 6E}). Moreover, the overwhelming majority of encoding models have negative slopes (meaning their performance decreases over time), while the average decoding model has a positive slope.

We next investigated the differences between the tuning functions for our single-unit encoding models. The vectors of weights of our models can be thought of as tuning functions that map behavioral representations to neural activity \citep{bib53,bib68,bib69}. We find that as the amount of time between training sets increases, the cosine distance of our tuning functions increases very gradually (weighted least squares, $p=0.0427$, \textbf{Figure 6F}). Thus, beyond performance loss, the temporal degradation in single-unit semantic encoding is accompanied by slight drift in the semantic tuning functions themselves. Therefore, while maximizing recording stability is an important avenue for the development of consistent semantic coding models, additional model-level corrections (such as aligning latent neural embeddings; \citealp{bib70}) are also likely needed to account for inherent neural variability across time.

\clearpage

\begin{figure}[H]
\centering
\includegraphics[width=\textwidth,height=0.86\textheight,keepaspectratio]{\figdir/Figure_6_Jun11.pdf}
\label{fig:figure6}
\end{figure}
\figlegend{\textbf{Figure 6 -- Functional drift.} \textbf{A.} Boxplot with line of best fit showing encoding model performance changes across time in days. \textbf{B.} Scatterplot of fraction of significant semantic encoding units as a function of test distance in days. \textbf{C.} Boxplot with line of best fit showing decoding model performance changes across time in days. \textbf{D.} Scatterplot of fraction of significant semantic decoding models as a function of test distance in days. \textbf{E.} Density plots showing distribution of encoding model performance slopes (green) and decoding model performance slopes (purple). \textbf{F.} Scatterplot with line of best fit showing changes in tuning function distance across time in hours.}

\clearpage

\section{Discussion}\label{sec10}

We developed a pipeline for processing of incidental speech data recorded 24/7 in the EMU. It records all speech and brain activity from implanted Behnke-Fried microwire bundles during the patient's stay (6+ days). Our pipeline includes automated spike sorting, speech transcription, and video-assisted diarization. We used this pipeline to generate reliable encoding and decoding models, whose performance we quantified by using ground-truth datasets consisting of manually annotated speech and manually sorted spikes. One of the most difficult practical, if not technical, problems with such pipelines is developing tools to assess their quality. Here we used a painstaking manual curation of naturalistic data involving both hand-sorting of spikes and precise manual alignment of all speech during typical sessions. This dataset allowed us to precisely quantify the accuracy of our much larger naturalistic dataset.

We found that both encoding and decoding models are robust and reliable, and decoding models trained on longitudinal incidental speech performed better than their controlled-task counterparts. These results indicate that the level of neural responding to language is such, and of such a type, that we can infer reliable, generalizable models from passively recorded speech. Moreover, they tell us that currently available methods are able to solve historically difficult problems in transcription, diarization, and spike sorting sufficiently well as to generate reliable models. This is important for applications in which bespoke data collection is impractical, such as in the continuous updating of brain-computer interfaces, or in scientific projects requiring large amounts of data, or long periods of time, or very infrequent words.

We found our encoding model performance scales with the size of the training dataset. Notably, even with large datasets, we did not find a saturation point for performance as a function of training set size. This result argues for the potential value of (large) naturalistic datasets in training models, especially when current approaches struggle to reach decoding accuracies useful for clinical application \citep{bib82,bib131,bib132,bib133,bib16,bib17}. Critically, this performance increase is achievable using automated pipelines, making the use of large datasets feasible; this is important because, as training sets grow larger, data preprocessing increasingly becomes a bottleneck.

Current approaches to modelling speech are often split between those that use single units and those that use local field potentials (LFPs). These two measures offer complementary, and partially overlapping, views of brain activity. Here we focus on spiking activity. While LFPs are relatively straightforward to analyze, and require little preprocessing, spiking activity requires detection (requiring distinguishing spikes from spike-like noise) and isolation (assigning spikes to the neurons that generate them), which is more difficult. Many automated and partially automated systems exist \citep{bib87,bib128,bib108,bib130}, but none approach the accuracy of careful manual annotation. Here, we show that perfect neuron-level spike isolation (as defined by that achievable by manual sorting) is not necessarily required for neural modeling, and current automated systems---specifically, spike autothresholding---are sufficient to generate reliable spike-based encoding and decoding models.

We found that we can recover significant language representations in all recorded brain regions, including the anterior cingulate cortex, hippocampus, and the centromedial head of the thalamus. These findings are inconsistent with the idea that language is strictly localized to specific language network regions, but are consistent with a growing body of neuroimaging work showing widespread brain representation of language \citep{bib2,bib8,bib140,bib139,bib3,bib138}. Our results therefore extend these ideas from neuroimaging to the domain of spiking activity. They also open the door to large data-driven studies of subtler differences between brain regions, including variation in semantic content and grammatical role.

Why did the large increase in longitudinal training data improve our population decoding models but not our channel-level encoding models? Our encoding models were fit to threshold-crossing activity from individual recording channels, rather than to spike-sorted single neurons; therefore, longitudinal nonstationarity could affect them in at least two ways. First, physical changes at the electrode--tissue interface can alter which neuron or mixture of neurons contributes to threshold crossings on a given channel. Thus, the same channel across days may not correspond to the same local neural source, making the apparent channel--stimulus relationship unstable even without any change in the underlying neural representation \citep{bib99}. Second, even in the absence of physical electrode drift, the neural representation itself may change over time. Such representational drift has been observed across sensorimotor and sensory systems \citep{bib98,bib141,bib142}. This is important because an encoding model asks what a particular channel's threshold-crossing activity is tuned to. Increasing the amount of training data should improve estimation by expanding the sampled stimulus space and providing more observations of similar semantic conditions, but only under the assumption that each channel's tuning function remains sufficiently stationary. If a channel's apparent tuning changes because the electrode samples a different local neuronal mixture, or because the same neurons undergo true representational drift, then pooling more longitudinal data can combine inconsistent channel--stimulus mappings rather than improve the model.

Population decoding models are expected to be more robust to this form of drift. Unlike single-unit encoding models, a decoder is not required to preserve a fixed relationship between a semantic feature and any one neuron. Instead, it can exploit redundancy across units, down-weight unstable or weakly informative units, and read out information from population dimensions that remain predictive over time. This interpretation is consistent with work in mouse visual cortex showing that individual neurons can change their activity rates and tuning over minutes to days, while the relationships among population activity patterns remain relatively stable and stereotyped \citep{bib142}. It is also consistent with evidence that sensory cortical population codes can become reliable through redundancy and structured inter-area co-fluctuations, making population decoding robust to day-to-day variability in individual cells \citep{bib143}. More generally, task-relevant population dynamics can be recovered without relying on stable spike sorting of individual units \citep{bib61}, and low-dimensional latent dynamics or manifolds can remain stable even when the recorded units themselves change, supporting reliable decoding over long timescales \citep{bib100,bib101}. Longitudinal speech-neuroprosthesis work further supports this population-level view, showing that decoders trained and adapted using cross-day data can maintain high accuracy over months of use \citep{bib102}. Thus, our results are consistent with an encoding--decoding asymmetry: single-unit encoding models directly expose drift in unit-level tuning, whereas population decoding models can average over, reweight, or align across unstable units to recover more stable population-level structure. This difference may explain why additional longitudinal data improved decoding performance but provided little benefit for naturalistic encoding models.

Variability in attention engagement is a major potential source of noise in incidental language studies. During dedicated tasks, patients attend closely; in the naturalistic speech settings, the patient decides what to attend to and when. It is highly unlikely that the patient is attending to the same degree at all times, or even attending at all most of the time. Neural activity likely covaries less with speech the patient is not attending to \citep{bib35,bib74,bib75}. Thus, one area for improvement in future automated models would be developing techniques for monitoring and controlling for attention levels, such as gaze-tracking or pupillometry \citep{bib121,bib122}. Future studies will need to address this shortcoming as well as others---including deficits in single-unit recording stability---in order to develop models from passively-recorded naturalistic data that are not only scientifically meaningful but also clinically usable.
\clearpage
\section{Methods}\label{sec11}

\subsection{Electrophysiology Recording}\label{subsec1}
All recorded epilepsy patients (8M, 13F, ages 18-54 years, average age: 34.5 $\pm$ 2.2 years) were undergoing neural monitoring for seizure localization at the Epilepsy Monitoring Unit (EMU) at Baylor St. Luke’s Medical Center in Houston, TX. Patients had Behnke-Fried intracranial electrodes placed in various regions of the brain that differed per patient, with a subset of depth electrodes including microwires from which single unit data was recorded and analyzed for this study.

\subsection{Behavioral Recording}\label{subsec2}
Two Blue Yeti microphones, one directly above the patient bed and another in the corner of the room, recorded all ambient audio in the EMU, which included frequent conversations between the patient and others. Only the audio from the microphone above the patient bed was used for our analysis as it tended to record patient speech with more clarity. Eight Flir cameras located at various positions throughout the room recorded video of patient bed from various angles, capturing patient motions, expressions, and activities. Camera quality was valididated via manual inspection per patient and the camera with the clearest view of the patient's face was used for analysis.

\subsection{Voice Activity Detection}\label{subsec3}
To identify sections of our recording containing human speech, we used silero\_vad (GitHub: \href{https://github.com/snakers4/silero-vad}{snakers4/silero-vad}, \citealp{bib43}), which is widely used in many speech processing pipelines and is considered state-of-the-art for voice activity detection, including in noisy, naturalistic environments. The model was provided audio from the room mic in ten minute chunks, and it outputted intervals of time predicted to contain human speech. Intervals that were less than one minute apart were merged up to a maximum of thirty minutes and considered continuous speech chunks. Each speech chunk was then preprocessed independently.

\subsection{Transcript Assembly}\label{subsec4}
To transcribe audio in the EMU, We ran each speech chunk through a modified version of OpenAI's Whisper called WhisperX \citep{bib39} that uses acoustic embeddings to improve word segmentation. We then ran all words through GPT-2 (\href{https://huggingface.co/openai-community/gpt2-large}{openai-community/gpt2-large}) with a maximum within-chunk context window of 200 words. We used the last layer of each word's hidden state to derive semantic embeddings of spoken speech in the EMU.

\subsection{Active Speaker Detection}\label{subsec5}
To identify sections of speech chunks with patient speech, we used an Active Speaker Detection (ASD) model called Lightweight and Robust Network for Active Speaker Detection (LR-ASD, GitHub: \href{https://github.com/Junhua-Liao/LR-ASD}{Junhua-Liao/LR-ASD}). This multimodal model is able to use both acoustic and facial features to predict whether a given person is speaking in a video frame. Since our video stream does not contain audio, we used an internal Python package (\href{https://github.com/BCM-Neurosurgery/video-sync-nbu}{BCM-Neurosurgery/video-sync-nbu}) to align each speech chunk with its associated video before providing it to LR-ASD. We then labeled segments of speech (as identified by WhisperX, roughly corresponding to a sentence) as patient speech if more than 50\% of the frames during that segment were predicted to contain the patient speaking by LR-ASD.

\subsection{Quality Metrics}\label{subsec6} 
For every sentence, we calculated its corresponding CTC log-likelihood, WhisperX quality score, spectral entropy, and DNSMOS quality scores. We identified outlier sentences using the interquartile range of each distribution and removed these sentences. This resulted in an average of 14.7\% ($\pm$ 1.5\%) of our sentences being considered low quality and filtered out before downstream analysis.

The connectionist temporal classification (CTC) loss quantifies the likelihood of a transcription given its associated audio segment. Using a pretrained wav2vec model, which is known to generate phonetically accurate acoustic embeddings of speech \citep{bib83}, we calculated embeddings of each audio segment and found its CTC loss value with its corresponding WhisperX transcript. We specifically used the \href{https://huggingface.co/facebook/wav2vec2-large-960h-lv60-self}{wav2vec2-large-960h-lv60-self} weights for the model. The CTC loss for transcript y and audio segment x is formalized below as:
\[
\mathcal{L}_{CTC}=-\log{p(y \mid x)}=-\sum_{\pi}\prod_{t=1}^{t}p(\pi_{t}\mid x)
\]
where $\pi_{t}$ is a symbol set in the set of all acoustic frame-level paths that lead to our target transcript. The precise formulation of frame-level paths as well as the algorithmic implementation of the loss can be found in the original Connectionist Temporal Classification paper \citep{bib84}. Our CTC log-likehood score is simply defined as the negative of the CTC loss value:
\[
score_\mathit{ctc}=-\mathcal{L}_\mathit{CTC}
\]

Spectral entropy measures the complexity and randomness of a signal and can serve as a proxy for assessing the noisiness of an audio segment (clean speech is typically concentrated in select frequency bands), with the aim being to exclude especially noisy segments. To compute spectral entropy, we transformed each audio segment into a power spectrogram time series with the short-time Fourier transform (STFT). We normalized each spectrogram across frequency bins to form a spectral probability distribution per time bin $P(f, t)$. We then computed the Shannon entropy of this distribution per time bin:
\[
H(t)=-\sum_{f}P(f,t)\log_{2}P(f,t)
\]
We took the average of $H(t)$ across time bins to calculate a segment-wide measure of spectral dispersion which we report as the spectral entropy score.

We also filtered audio segments via quality scores provided directly by the WhisperX pipeline, which quantify the pipeline’s alignment confidence between audio and its associated transcript. We lastly filtered using scores provided by DNSMOS, a neural estimator of perceptual speech quality trained to approximate human-labeled quality scores on measures of speech clarity and background noise. We refer readers to the respective papers for specifics on how these scores are calculated \citep{bib39,bib85}.

Together these filtration metrics provide a means to exclude low-quality transcripts that would negatively impact our downstream modeling performance without resorting to extensive manual checks of our dataset. We aimed to capture measures across orthogonal modes of contamination, especially those concerning raw audio quality and transcript alignment probability.  

There was no resulting significant difference in encoding model performance between the filtered and unfiltered dataset (Welch's t-test, p=0.431), but all analyses were performed on the filtered dataset to maximize the quality of the behavioral dataset.

\subsection{Spike Auto-thresholding}\label{subsec7}
To capture spike counts across channels per word, we used the following bisection autothresholding algorithm:

Firstly, raw extracellular voltage traces were bandpass filtered to spike band range (300-6000 Hz) using a zero-phase forwards and backwards pass of the 4th order Butterworth filter. 

For each channel we calculated the median absolute deviance (MAD) of the bandpassed voltage signal as the noise scale for calculating the spike threshold:
\[
\sigma_{c}=1.4826\cdot \mathrm{median}(|x_{c}^{(f)}-\mathrm{median}(x_{c}^{(f)})|)
\]
Spikes were detected as local minima in the signal below $-k\sigma_{c}$, where $k$ is solved via bisection. 

To prevent double-counting of spikes an absolute refractory period of 1 ms was required for the detection of a subsequent spike. 
To solve $k$ we calculated the global firing rate across channels for every bisection iteration of $k$ and iterated until the global firing rate was within 0.5 Hz of 20 Hz or the maximum number of iterations had been exceeded (25). The global firing rate of $k$ is defined as:
\[
r(k)=\frac{1}{CT}\sum_{c=1}^{C}N_{c}(k)
\]
where $C$ is total channels, $T$ is recording duration, and $N_{c}(k)$  is the recorded number of spikes for a particular channel $c$ and constant $k$

Detected spike times were converted to a binary spike train per channel sampled at 1 kHz where the value was set to 1 for each 1 ms time bin where any number of spikes occurred. Subsequently, spike counts per word was calculated as the number of spikes occurring during the total duration of the word.
.
\subsection{Poisson GLM}\label{subsec8}
Our semantic encoding model involved using GPT-2 embeddings to predict firing rate. Our model inputs are defined below for every $i^{th}$ input:
\begin{itemize}[leftmargin=5em,label=\textendash,itemsep=3pt,topsep=4pt]
  \item feature vector $x_{i}\in \mathbb{R}^d$, where $d$ is the dimensionality of the principal components (PCs) of the last GPT-2 hidden state (1280 → 100). PCs were used to maximize information across dimensions and reduce model complexity.
  \item spike count vector $y_{i}\in \mathbb{R}^K$, where $K$ is the number of channels recorded from and each value is the associated spike count for word $i$
  \item word duration $\Delta_{i}$ in milliseconds
  \item offset $o_i\log(\Delta_i)$
\end{itemize}

Our GLM predicts the associated spike count of an input vector with the log link function:
\[
\begin{aligned}
\log(\mu_{ik}) &= x_i^T w_k + b_k + o_i, \quad \text{implying} \\
\mu_{ik} &= \Delta_i \exp(x_i^T w_k + b_k)
\end{aligned}
\]

Parameters were estimated by minimizing the summed Poisson negative log-likelihood across observations and neurons with an L2 (ridge) penalty. All $K$ neurons were fit jointly using matrix operations in PyTorch, yielding a single scalar loss optimized with the L-BFGS algorithm.
\begin{gather*}
\mathcal{L}(W,b)=\sum_{i=1}^N\sum_{k=1}^K(\mu_{ik}-y_{ik}\log(\mu_{ik}))+\frac{1}{2}\sum_{i=1}^{N}\sum_{k=1}^{K}\alpha_k\lVert w_{ik}\rVert^2
\end{gather*}
We used 5-fold outer cross-validation to evaluate held-out performance. Within each outer training split, we selected k via 5-fold inner cross-validation by maximizing mean validation log-likelihood per neuron over a grid of candidate $\alpha$ values. To accelerate hyperparameter search, we warm-started optimization across successive $\alpha$ values (evaluated from high to low regularization) by initializing each fit with the solution from the previous $\alpha$ on the same fold. Values for $\alpha$ were selected from 30 candidates log-spaced in the range $[10^{-3}, 10^3]$.

\subsection{Pseudo R\textsuperscript{2}}\label{subsec9}
Pseudo R\textsuperscript{2} values were calculated as the average Pseudo R\textsuperscript{2} value across 5 held out test sets in a 5-fold cross validated scheme. Pseudo R\textsuperscript{2} definition is below:
\begin{gather*}
\psi R^2=1-\left(\frac{\log\mathcal{L}_\mathit{real}}{\log\mathcal{L}_\mathit{null}}\right), \quad\text{where}\\
\log\mathcal{L}=\sum_{i=1}^{N}y_i\log\mu_i-\mu_i\log(|\Gamma(y_i+1)|), \quad\text{where $\Gamma$ is the gamma function}
\end{gather*}
$\mu_{ik,\mathit{real}}$ is the individual prediction from our Poisson GLM, whereas $\mu_{ik,\mathit{null}}$ is the prediction from our null model predicting spike counts solely from average firing rate of each channel, formalized as:
\begin{gather*}
\mu_{ik,\mathit{null}}=\bar{\lambda}_k\Delta_i, \quad \text{where}\\
\bar{\lambda}_k=\frac{\sum_i^{N_\mathit{train}}y_{ik}}{\sum_i^{N_\mathit{train}}\Delta_i}, \quad\text{and}\\
\text{$y_{ik}$ is the number of spikes in word $i$ for channel $k$ and $\Delta_i$ is the duration of word $i$}
\end{gather*}

\subsection{XGBoost Decoder}\label{subsec10}
To decode semantic category labels from population firing rate activity, we trained multiclass gradient-boosted decision tree classifiers using XGBoost. For each word $i$, the feature vector $x_i$ is defined as the average firing rate across the word’s duration for all $K$ recorded channels. Per-bin firing rate was calculated by convolving a gaussian kernel with 50 ms standard deviation across our autothresholded binary spike train. The average firing rate was calculated as the mean of all firing rates whose bins overlapped with the word duration. The output label $y_i$ is the corresponding semantic category of the word, which is classified from the word2vec embedding of the word (see: \textbf{4.10}). The model was optimized using the multiclass log loss:
\[
\mathcal{L}=-\frac{1}{N}\sum_{i=1}^N\sum_{c=1}^C\mathbf{1}[y_i=c]\log\hat{p}_{ic},\quad\text{where $C$ is the number of classes}
\]
Hyperparameters including max\_depth, min\_child\_weight, gamma, subsample, colsample\_bytree, reg\_lambda, reg\_alpha, learning\_rate, and n\_estimators were tuned using 5-fold Randomized Search Cross Validation (implemented in \href{https://scikit-learn.org}{scikit\_learn}). The highest performing hyperparameters were then used to train a model on the full label-stratified training set and evaluated on a single held-out test set.

To reduce prevalence bias, we excluded function words and downsampled remaining classes to the median class count via random sampling. To prevent sampling bias, we fit the XGBoost model with the above methodology for 20 resamples. To reduce compute time, hyperparameter tuning was only performed for the first 3 resamples and the remaining resamples used the hyperparameters from the initial 3 fits that maximally reduced model complexity.

Scoring for cross-validation of our semantic XGBoost decoders was evaluated using the F1 macro score:
\begin{gather*}
F_1=\frac{1}{C}\sum_{c=1}^CF_{1,c}, \quad\text{where}\\
F_{1,c}=\frac{2TP}{2TP+FP+FN}, \quad\text{for class $c$}
\end{gather*}

Model performance was reported as the general accuracy score of the model on the held-out test set:
\[
acc=\frac{TP+TN}{TP+TN+FP+FN}
\]

\subsection{Semantic Category Classifier}\label{subsec11}
Semantic categories were determined from an XGBoost classifier trained to predict class labels from a word’s word2vec embedding. Initial semantic labels were provided by the conversation dataset through clustering methods described in \citet{bib53}. To extrapolate these labels to other datasets, we generated word2vec embeddings of each conversation word with an associated label via fastText word vectors pretrained on the English sections of Common Crawl and Wikipedia (using the model downloadable at \href{https://fasttext.cc/docs/en/crawl-vectors.html}{fasttext.cc}. Then, we trained an XGBoost classifier with the same randomized hyperparameter search described in \textbf{4.10} on a stratified training set. The $F_1$ macro score on the held-out test set was 0.96, indicating the semantic categories described in \citet{bib53} are reliably recoverable from their word2vec embeddings. 

\subsection{Task Information}\label{subsec12}
Two control tasks, colloquially referred to as “podcast” and “conversation”, were used to compare the relative performance gain of naturalistic speech-trained models.

The podcast task involved patients listening to a 47-minute sequence of six autobiographical stories as part of The Moth Radio Hour, totalling 7,346 words of natural speech. They were asked to pay full attention throughout the duration of the task session

The conversation task involved patients engaging in casual conversations with research staff. These conversations varied in time per patient but were usually between 30-60 minutes and length and covered a range of topics at the patient’s discretion. Conversations frequently included patient visitors as well if present. 

\subsection{Sentence Matching Algorithm}\label{subsec13}
To evaluate whether a particular sentence in our corrected dataset had a match in the naturalistic speech dataset, we searched all sentences in a 60 second window around the query sentence. For all sentences with a fuzzy token set ratio $\ge$ 60 (implemented as token\_set\_ratio in the \href{https://rapidfuzz.github.io/RapidFuzz/}{RapidFuzz} Python package), we aligned the tokens from the corrected sentence to the naturalistic sentence to calculate the Word Error Rate (WER) between the sentences. This was done via a 2D dynamic programming algorithm that calculates the minimum amount of token substitutions (S), deletions (D), and insertions (I) needed to transform the reference sentence into the target sentence (token-level Levenshtein edit distance). A token was considered a match (M) between two sentences (i.e., edit distance of 0) based on their character-level Levenshtein distance as below for tokens A, B:
\[
M = \begin{cases} 
      length(A)\le4&\mathbb{I}[Ldist(A,B)\le1] \\
      length(B)>4&\mathbb{I} [Ldist(A,B)\le2]
   \end{cases}
\]
WER is formalized for a reference sentence as:
\[
WER=\frac{S+D+I}{N},\quad\text{where $N$ is the number of reference tokens}
\]
And WER similarity as:
\[
SIM_\mathit{WER}=1-WER
\]
The sentence with the highest $SIM_\mathit{WER}$ was taken as the corrected sentences’ match. If no sentence had $SIM_\mathit{WER} >0.5$ the sentence was marked unmatched.

\subsection{Sentence Match Summary Statistics}\label{subsec14}
The token-level Levenshtein distance described in \textbf{4.13} calculates the matched words between the reference and naturalistic sentence. To define the distribution of word timestamp errors, we simply calculated the median timestamp of each word referenced to the start of its sentence and calculated the difference between the reference and naturalistic timestamp. The timestamp error for word pair $i$ in sentence pair $j$ is defined as:
\begin{gather*}
\epsilon_i=m_{ai}-m_{ri},\quad\text{where}\\
m_{ai}=\frac{start_{ai}+end_{ai}}{dur_{ai}}-start_{aj},\quad\text{and}\\
m_{ri}=\frac{start_{ri}+end_{ri}}{dur_{ri}}-start_{rj}
\end{gather*}

The same error metric was calculated for each sentence $i$ as well to compare sentence-level temporal errors. The resulting distributions describe the extent of temporal errors in sentence-level segmentation between corrected and naturalistic sentences.

We defer to the original Sentence-BERT paper for a detailed explanation of the model used to transform sentences into semantic embeddings suitable for cosine distance comparison \citep{bib54}.

\subsection{Permutation Testing}\label{subsec15}

To evaluate whether each neuron’s GLM performance was significant compared to predictions based on word duration or other latent structural features, we randomly permuted the GPT embeddings 50 times for each cross fold and refit our Poisson GLM with the same vector. This destroyed the semantic relationship between input and output whilst maintaining word duration as labels were not permuted instead. PCA was recalculated for the newly permuted training set per fold. The p-value per fold is calculated by observing the proportion of times the log-likelihood \textbf{4.8} value of our permuted model is greater than or equal to our real model:
\[
p=\frac{1+\sum_{n=1}^{N_\mathit{shuf}}\mathbf{1}[\log\mathcal{L}_{\mathit{shuf},n}\ge\log\mathcal{L}_\mathit{real}]}{N_\mathit{shuf}+1}
\]
The p-value across folds for a particular neuron is calculated the same way except by using the average $log\mathcal{L}$  values across folds and per permutation:
\begin{gather*}
\log\mathcal{L}_\mathit{real}=\frac{1}{K}\sum_{k=1}^K\log\mathcal{L}_{\mathit{real},k,n}\\
\log\mathcal{L}_{\mathit{shuf},n}=\frac{1}{K}\sum_{k=1}^K\log\mathcal{L}_{\mathit{shuf},k,n}\\
p=\frac{1+\sum_{n=1}^{N_\mathit{shuf}}\mathbf{1}[\log\mathcal{L}_{\mathit{shuf},n}\ge\log\mathcal{L}_\mathit{real}]}{N_\mathit{shuf}+1}
\end{gather*}

\subsection{Spike Sorting}\label{subsec16}
To evaluate the effect of neural preprocessing techniques on our downstream , we also sorted raw voltage traces with the algorithm included as part of the wave\_clus MATLAB package \citep{bib87}. The resulting sorted clusters were then manually evaluated and filtered were then manually evaluated and filtered with the wave\_clus GUI. Non-spike noise was removed and each signal was classified as multi or single unit using several criteria: consistent spike waveforms, waveform shape, exponentially decaying ISI histogram, and refractory period (1 ms) violations below 5\%. Sorted signals were included for both single and multi unit activity.

\subsection{Model Performance Comparisons}\label{subsec17}
All statistical analyses were performed using two-sided tests unless otherwise specified, with significance assessed at $\alpha = 0.05$.

Word error rate similarity and SBERT cosine similarity were compared across matched sentences using a paired Wilcoxon signed-rank test. Each matched sentence contributed one pair of observations, denoted
\[
(w_i, c_i),
\]
where $w_i$ is the WER similarity and $c_i$ is the cosine similarity for sentence $i$. Differences were computed as
\[
d_i = c_i - w_i.
\]
The Wilcoxon signed-rank statistic was computed after excluding zero differences:
\[
W = \sum_{i: d_i > 0} R_i,
\]
where $R_i$ is the rank of $|d_i|$ among nonzero paired differences. The null hypothesis was
\[
H_0: \operatorname{median}(d_i) = 0.
\]
Both two-sided tests and one-sided tests for
\[
H_A: \operatorname{median}(d_i) > 0
\]
were computed in the code.

For model condition comparisons (all-speech vs. patient-only, whole-stay vs. per-day, etc.), group differences were assessed using Welch's two-sample $t$-test. The statistical tests were performed on the pooled unit-level values for encoding analyses and pooled resample-level values for decoding analyses.

For two groups with means $\bar{x}_1$, $\bar{x}_2$, variances $s_1^2$, $s_2^2$, and sample sizes $n_1$, $n_2$, the test statistic was
\[
t =
\frac{\bar{x}_1 - \bar{x}_2}
{\sqrt{s_1^2/n_1 + s_2^2/n_2}}.
\]
Degrees of freedom were estimated using the Welch--Satterthwaite approximation:
\[
\nu =
\frac{
\left(s_1^2/n_1 + s_2^2/n_2\right)^2
}
{
\frac{(s_1^2/n_1)^2}{n_1 - 1}
+
\frac{(s_2^2/n_2)^2}{n_2 - 1}
}.
\]

\subsection{Linear Modeling}\label{subsec18}

Weighted least-squares regression was used for relationships between model performance and dataset or drift variables. For observations $(x_i, y_i)$, the fitted model was
\[
y_i = \beta_0 + \beta_1 x_i + \varepsilon_i.
\]
The WLS estimate minimized
\[
\sum_{i=1}^{n} w_i
\left(y_i - \beta_0 - \beta_1 x_i\right)^2,
\]
with solution
\[
\hat{\boldsymbol{\beta}}
=
(X^\top W X)^{-1}X^\top W\mathbf{y}.
\]

For patient-speech fraction vs. epoch model performannces (\textbf{Figure 3E}, \textbf{Figure 4F}), $x_i$ was the patient-speech fraction within an epoch and $y_i$ was encoding pseudo R\textsuperscript{{2}} or decoding accuracy, respectively. WLS fits were performed separately for each patient. Weights were assigned by inverse bin density of the predictor, using Freedman--Diaconis bins capped at 20 bins, so overrepresented speech-fraction ranges contributed less strongly to the fit.

For populution size vs. decoding performance (\textbf{Figure 4E}), $x_i$ was population size and $y_i$ was decoding accuracy. The same inverse-density WLS procedure was used.

For the above regressions, slope significance was assessed using a two-tailed permutation test with 500 permutations. The observed slope $\hat{\beta}_1$ was compared with slopes obtained after shuffling predictor values:
\[
p =
\frac{
1 + \sum_{b=1}^{B}
\mathbf{1}
\left(
|\hat{\beta}_{1}^{(b)}| \geq |\hat{\beta}_{1}|
\right)
}
{B + 1},
\qquad B = 500.
\]

To assess functional drift (\textbf{Figure 6A}, \textbf{Figure 6C}), WLS regression was used to test whether the model performance changed with days since training. Each unique day-offset group was weighted equally:
\[
w_i = \frac{1}{n_{g(i)}},
\]
where $n_{g(i)}$ is the number of observations in the day-offset group containing observation $i$.

Slope significance was assessed using a WLS $t$-test:
\[
t =
\frac{\hat{\beta}_1}{\operatorname{SE}(\hat{\beta}_1)}.
\]
where \(\hat{\beta}_1\) is the estimated WLS slope and
\(\operatorname{SE}(\hat{\beta}_1)\) is its standard error.

For encoding vs. decoding drift comparison (\textbf{Figure 6E}), drift slopes were computed separately for each patient and training day after min--max normalizing performance within each training-day series.

\subsection{Epoch Analyses}\label{subsec19}
To assess short-term changes in encoder/decoder model performance, we fit models on shorter timescales than our full-stay or per-day analyses. These were trained and tested identically to models in \textbf{4.8}, except each model was trained on a particular 2-hour block from 9 AM - 11 PM (local time).

Prior to weighted-least-squares regression on epoch-level pseudo R\textsuperscript{{2}} values, we removed unit--epoch observations whose pseudo R\textsuperscript{{2}} magnitude fell in the catastrophic numerical-failure tail using a data-derived elbow criterion. Let $\{a_j\} = \{|\psi R^2_j|\}$ denote the set of absolute pseudo R\textsuperscript{{2}} values across all finite unit--epoch fits. We evaluated tail quantiles on a fine grid $s \in [s_0, 100]$ (with $s_0 = 95$, $n_\text{grid} = 501$ points), and transformed to a log scale:
\[
q_s = \mathrm{Quantile}(\{a_j\},\, s), \qquad
\ell_s = \log_{10}(q_s + 1).
\]
Both axes were then normalized to $[0,1]$:
\[
\tilde{s} = \frac{s - s_{\min}}{s_{\max} - s_{\min}}, \qquad
\tilde{\ell}_s = \frac{\ell_s - \ell_{\min}}{\ell_{\max} - \ell_{\min}}.
\]
The elbow was located at the quantile level that maximized the vertical distance below the diagonal of this normalized curve (Kneedle criterion):
\[
s^{*} = \arg\max_{s}\!\left(\tilde{s} - \tilde{\ell}_{s}\right),
\]
and the corresponding fence value $\tau = q_{s^{*}}$ defined the inclusion threshold. Unit--epoch rows with $|\psi R^2| > \tau$ were excluded from WLS analyses; values within the fence but outside the display range $[-1, 1]$ were retained in the regression while the plot axis remained fixed.

\subsection{Variance Comparison}\label{subsec20}
We compared the relative dispersion of encoding and decoding model performance across recording epochs using the coefficient of variation (CV). Encoding dispersion was computed from all positive, significant single-unit pseudo R\textsuperscript{{2}} values pooled across patients and epochs. Decoding dispersion was computed from all positive accuracy values from individual decoding resamples pooled across patients and epochs.

For a set of performance values
\[
x = \{x_1, x_2, \ldots, x_n\},
\]
the coefficient of variation was defined as
\[
\mathrm{CV}
=
\frac{s}{\bar{x}},
\]
where
\[
\bar{x}
=
\frac{1}{n}\sum_{i=1}^{n} x_i
\]
is the sample mean and
\[
s
=
\sqrt{
\frac{1}{n-1}
\sum_{i=1}^{n}
(x_i - \bar{x})^2
}
\]
 is the sample standard deviation. Thus, for encoding and decoding,
\[
\mathrm{CV}_{\mathrm{enc}}
=
\frac{s_{\mathrm{enc}}}{\bar{x}_{\mathrm{enc}}},
\qquad
\mathrm{CV}_{\mathrm{dec}}
=
\frac{s_{\mathrm{dec}}}{\bar{x}_{\mathrm{dec}}}.
\]

Only positive values were included in the CV calculation, because CV is defined relative to the mean and is most interpretable for positive-valued measurements.

To test whether encoding and decoding differed in relative dispersion, we used the Feltz--Miller asymptotic test for equality of two coefficients of variation. The null hypothesis was
\[
H_0:
\mathrm{CV}_{\mathrm{enc}}
=
\mathrm{CV}_{\mathrm{dec}},
\]
with alternative hypothesis
\[
H_A:
\mathrm{CV}_{\mathrm{enc}}
\neq
\mathrm{CV}_{\mathrm{dec}}.
\]

Let \(n_1\) and \(n_2\) denote the number of encoding and decoding observations, and let
\[
c_1 = \mathrm{CV}_{\mathrm{enc}},
\qquad
c_2 = \mathrm{CV}_{\mathrm{dec}}.
\]
A pooled CV estimate was computed as
\[
c_p =
\frac{(n_1 - 1)c_1 + (n_2 - 1)c_2}
{n_1 + n_2 - 2}.
\]
The Feltz--Miller test statistic was then computed as
\[
T =
\frac{
(n_1 - 1)(\log c_1)^2
+
(n_2 - 1)(\log c_2)^2
-
\frac{
\left[
(n_1 - 1)\log c_1
+
(n_2 - 1)\log c_2
\right]^2
}
{n_1 + n_2 - 2}
}
{
0.5 + c_p^2
}.
\]
Under the null hypothesis, this statistic was evaluated against a chi-squared distribution with one degree of freedom:
\[
T \sim \chi^2_1.
\]

\subsection{Region Decoder Sampling}\label{subsec21}
To assess the relative performance of brain regions in semantic category decoding, we fit population decoders using only units from specific regions. Those units that were on the border of their intended target or otherwise missed it according to visual inspection of the post-operative, coregistered CT scan in RAVE \citep{bib88} were excluded from this analysis. To remove unit count as a confounding factor, we fit each population decoder with only 8 units. If we had more than 8 units for a particular region, we sampled 8 units 5 times and fit 10 category-sampled models for a total of 50 fits per brain region. Otherwise, we only fit 10 category-sampled models for patient regions where there were only 8 units. To increase training efficiency, we only performed full hyperparameter search on the first three sample models per region and then picked the hyperparameters that most reduced model complexity for subsequent fits. 

\subsection{Decimation Analysis}\label{subsec22}
To evaluate the performance growth of our encoding models, we fit Poisson GLMs from \textbf{4.8} using randomly sampled proportions of the per-day naturalistic dataset for each patient. We fit models using 1\%, 5\%, 10\%, 20\%, 30\%, 40\%, 50\%, 60\%, 70\%, 80\%, 90\%, and 100\% of the data with 100, 20, 10, 8, 7, 6, 5, 4, 3, 2, 2, and 1 resamples, respectively. 

\subsection{Hill Curve Fitting}\label{subsec23}
To evaluate whether our performance growth for encoding models followed a logistic growth curve, we fit a predictive model of pseudo R\textsuperscript{2} as a function of data proportion/number of words used to train the model using the Hill equation:
\[
\psi R^2=\frac{Ax^h}{K^h+x^h}+b, \quad \text{where $A$, $h$, $K$ are weights and $b$ is the bias term}
\]
The model was fit using L-BFGS with a weighted mean squared error loss (so as to equally consider proportions  regardless of sampling frequency) defined below:
\[
\mathcal{L}=\sum_{i=1}^Nw_i'(y_i-\hat{y}_i)^2, \quad \text{where $\hat{y}_i$ is the predicted $\psi R^2$, $y_i$ is the true $\psi R^2$, and $w_i'$ is the normalized weight}
\]
The weight term is defined as below:
\[
w_i=\frac{1}{\mathit{counts}(\mathit{bin}_i)}, \quad \text{where $\mathit{bin}_i$ is the proportion or histogram bin $i$}
\]
We then normalize the weight term for model fitting as below:
\[
w_i'=\frac{w_i}{\bar{w}}
\]
We then use the same permutation testing scheme as \textbf{4.15} to evaluate the significance of our model compared to chance. We also assessed the relative performance of the Hill model to a simple linear equation with a single weight and bias by calculating the differences in their loss value, negative values indicating our Hill curve achieved a better fit than a linear equation:
\[
\Delta\mathcal{L}=\mathcal{L}_\mathit{Hill}-\mathcal{L}_\mathit{linear}
\]

\subsection{Functional Drift}\label{subsec24}
To assess how our models performed on future data they were not trained on, we performed a functional drift analysis. To do this, we fit models identically to \textbf{4.8} except training data was restricted to the hours of 9 AM and 11 PM on a particular day and then we calculated test performance statistics on each day independently after the initial training day from the hours of 9 AM to 11 PM.

We then fit WLS regressions similar to \textbf{4.18} on the model test performance data to assess if model performance changed significantly over time.

\subsection{Tuning Function Drift}\label{subsec25}

To assess representational drift of our semantic encoding model, we fit independent Poisson GLMs for each 2-hour epoch of the patient stay similar to \textbf{4.19}, except we used a global PCA to transform all the data along the same basis to ensure weight distances reflected genuine differences in tuning rather than a change of basis.

To quantify representational drift, we compared the linear vector of weights for each model per-neuron using the cosine similarity score:
\[
\mathit{SIM}_\mathit{cos}=\frac{A\cdot B}{\lVert A\rVert \lVert B\rVert}
\]
To assess if weight distance changed as a function of time apart, we fit a WLS regression similar \textbf{4.18} of weight distance of models vs. time apart, binned by days apart in terms of calendar date (i.e., 9 PM on Tuesday and 11 AM on Wednesday would be considered one day apart in this scheme).

\backmatter

\bmhead{Ethics approval and consent to participate}
All experimental protocols were approved by the Baylor College of Medicine Institutional Review Board (IRB). All methods were carried out in accordance with the relevant guidelines and regulations. Informed consent was obtained from all participants and/or their legal guardians.

\bmhead{Code availability}
All analysis and processing scripts can be found at \url{https://github.com/BCM-Neurosurgery/speech_247_scripts}.

\bmhead{Use of generative AI}
Generative AI tools were used to assist with code refinement, troubleshooting of analysis and processing scripts, and copy editing of author-written manuscript text for clarity and readability. All analyses, interpretation, manuscript content, and final wording were reviewed and approved by the authors, who take full responsibility for the final manuscript.

%%===========================================================================================%%
%% If you are submitting to one of the Nature Portfolio journals, using the eJP submission   %%
%% system, please include the references within the manuscript file itself. You may do this  %%
%% by copying the reference list from your .bbl file, paste it into the main manuscript .tex %%
%% file, and delete the associated \verb+\bibliography+ commands.                            %%
%%===========================================================================================%%

\clearpage

% Align the reference list's left edge with the body text (the bibliography
% \list otherwise resets to the raw page margin, ~0.33in left of body text).
\begin{adjustwidth}{0.33in}{0pt}
\bibliography{sn-bibliography}% common bib file
\end{adjustwidth}
%% if required, the content of .bbl file can be included here once bbl is generated
%%\input sn-article.bbl

\end{document}
